\documentclass[11pt,a4paper]{article}
\usepackage{geometry}
\usepackage{tikz}
\usepackage{enumitem}
\usepackage{multirow}
\usepackage{pgfplots}
\usepackage[framemethod=tikz]{mdframed}%membuat frame

\geometry{lmargin=1.5cm,tmargin=1cm,rmargin=1cm,bmargin=1cm}
\renewcommand*\rmdefault{cmss}

\begin{document}

\begin{flushleft}

%-----------------------------------------------------Latihan
\begin{tikzpicture}
   \draw [draw=green,fill=yellow!30!white,rounded corners=5pt](0,0) rectangle (1.8,0.8);
   \node at (0.9,0.4){\large Latihan (File Ini Diedit dengan Commit Github};
\end{tikzpicture}\\
Selesaikanlah soal-soal berikut!\\[0.8em]
\begingroup
\renewcommand{\arraystretch}{1.6} % Default value: 1
\begin{tabular}{p{0.3cm}p{0.3cm}p{16cm}} 
   \bf{1.}&\multicolumn{2}{p{16cm}}{Tentukanlah penyelesaian dari PLSV berikut tanpa menulis lengkap langkah-langkahnya.}\\
    &a.&$4x+8=2x$\\
    &b.&$8y-14=5y-2$\\
    &c.&$\frac{3}{4}y=9$\\
    &d.&$\frac{x}{3}-1=3+\frac{x}{2}$\\
   \bf{2.}&\multicolumn{2}{p{16cm}}{Selesaikanlah PLSV berikut.}\\
    &a.&$3(4-2x)=3x-6$\\
    &b.&$4(2a-1)+10=6(4+a)$\\
    &c.&$\frac{1}{3}(p+2)=\frac{1}{4}(p-1)$\\
    &d.&$\frac{4x+2}{3}+\frac{2x+1}{2}=\frac{x+3}{4}$\\
   \bf{3.}&\multicolumn{2}{p{16cm}}{Tentukanlah penyelesaian dari PLSV berikut.}\\
    &a.&$4x-5=3x+9$\\
    &b.&$6a-5=4a-15$\\
    &c.&$-\frac{3}{4}x=6$\\
    &d.&$\frac{2}{3}x-15=3$\\
   \bf{4.}&\multicolumn{2}{p{16cm}}{Tentukanlah penyelesaian dari persamaan berikut.}\\
    &a.&$3x-10=2-x$\\
    &b.&$2(3p-6)=3p-9$\\
    &c.&$3(2y+1)-4(2y-3)=5$\\
    &d.&$\frac{1}{2}(x+3)=\frac{1}{4}(3x-1)$\\
   \bf{5.}&\multicolumn{2}{p{16cm}}{Carilah penyelesaian dari PLSV berikut.}\\
    &a.&$\frac{2}{3}(5x+2)=2(x+2)$\\
    &b.&$\frac{2y+4}{4}-\frac{4y-1}{3}=1$\\
    &c.&$\frac{2x-1}{4}+\frac{6x+4}{2}=4x+3$\\
   \bf{6.}&\multicolumn{2}{p{16cm}}{Tiga bilangan ganjil berurutan jumlahnya 51. Berapakah bilangan yang terkecil?}\\
   \bf{7.}&\multicolumn{2}{p{16cm}}{Suatu segitiga samakaki memiliki panjang sisi $x$ cm, $x$ cm, dan $x-5$ cm. Hitunglah keliling segitiga tersebut.}\\
   \bf{8.}&\multicolumn{2}{p{16cm}}{Berat $\frac{2}{5}$ karung beras adalah 8 kg. Berapa berat 1 karung beras?}\\
\end{tabular}
\endgroup\\[0.7em]

{\bf Kunci Jawaban}:\\
{\bf 1.} (a.) $x=-4$; \quad (b.) $y=4$ \quad  (c.) $y=12$ \quad (d.) $x=-24$\\
{\bf 2.} (a.) $x=2$; \quad (b.) $a=9$ \quad  (c.) $p=-11$ \quad (d.) $x=-\frac {1}{5}$\\
{\bf 3.} (a.) $x=14$; \quad (b.) $a=-5$ \quad  (c.) $x=-8$ \quad (d.) $x=27$\\
{\bf 4.} (a.) $x=3$; \quad (b.) $p=1$ \quad  (c.) $y=5$ \quad (d.) $x=7$\\
{\bf 5.} (a.) $x=2$; \quad (b.) $y=\frac {2}{5}$ \quad  (c.) $x=-\frac{5}{2}$\\

\newpage
{\bf Penyelesaian Lengkap}\\[0.3em]
\begin{minipage}{5.5cm}
{\color{blue} --------------------}\\
{\bf (1a)}\\
$4x+8=2x$\\
$4x-2x=-8$\\
$2x=-8$\\
$x=\frac{-8}{2}=-4$\\
{\color{blue} --------------------}\\
{\bf (1b)}\\
$8y-14=5y-2$\\
$8y-5y=-2+14$\\
$3y=12$\\
$y=\frac{12}{3}=4$\\
{\color{blue} --------------------}\\
{\bf (1c)}\\
$\frac{3}{4}y=9$\\
$(\frac{3y}{4}=9)\times 4$\\
$3y=9\times 4$\\
$3y=36$\\
$y=\frac{36}{3}=12$\\
{\color{blue} --------------------}\\
{\bf (1d)}--bonus--\\
$\frac{x}{3}-1=3+\frac{x}{2}$\\
$(\frac{x}{3}-1=3+\frac{x}{2})\times 6$\\
$x\cdot2 - 6=18 + x\cdot3$\\
$2x-6=18+3x$\\
$2x-3x=18+6$\\
$-x=24$\\
$x=\frac{24}{-1}=-24$\\
{\color{blue} --------------------}\\
{\bf (2a)}\\
$3(4-2x)=3x-6$\\
$12-6x=3x-6$\\
$-6x-3x=-6-12$\\
$-9x=-18$\\
$x=\frac{-18}{-9}=2$\\
{\color{blue} --------------------}\\
{\bf (2b)}\\
$4(2a-1)+10=6(4+a)$\\
$8a-4+10=24+6a$\\
$8a+6=6a+24$\\
$8a-6a=24-6$\\
$2a=18$\\
$a=\frac{18}{2}=9$\\
{\color{blue} --------------------}\\
{\bf (2c)}\\
$\frac{1}{3}(p+2)=\frac{1}{4}(p-1)$\\
$\left(\frac{1}{3}(p+2)=\frac{1}{4}(p-1)\right)\times 12$\\
$4(p+2)=3(p-1)$\\
$4p+8=3p-3$\\
$4p-3p=-3-8$\\
$p=-11$\\
{\color{blue} --------------------}\\
\\
\\
\\
\\
\end{minipage}
\begin{minipage}{6cm}
{\color{blue} --------------------}\\
{\bf (2d)}--bonus--\\
$\frac{4x+2}{3}+\frac{2x+1}{2}=\frac{x+3}{4}$\\
$\left(\frac{4x+2}{3}+\frac{2x+1}{2}=\frac{x+3}{4}\right)\times 12$\\
$4(4x+2)+6(2x+1)=3(x+3)$\\
$16x+8+12x+6=3x+9$\\
$28x-3x=9-14$\\
$25x=-5$\\
$x=\frac{-5}{25}=-\frac{1}{5}$\\
{\color{blue} --------------------}\\
{\bf (3a)}\\
$4x-5=3x+9$\\
$4x-3x=9+5$\\
$x=14$\\
{\color{blue} --------------------}\\
{\bf (3b)}\\
$6a-5=4a-15$\\
$6a-4a=-15+5$\\
$2a=-10$\\
$a=\frac{-10}{2}=-5$\\
{\color{blue} --------------------}\\
{\bf (3c)}\\
$-\frac{3}{4}x=6$\\
$(\frac{-3x}{4}=6)\times 4$\\
$-3x=9\times 4$\\
$-3x=24$\\
$x=\frac{24}{-3}=-8$\\
{\color{blue} --------------------}\\
{\bf (3d)}--bonus--\\
$\frac{2}{3}x-15=3$\\
$(\frac{2x}{3}-15=3)\times 3$\\
$2x-45=9$\\
$2x=54$\\
$x=\frac{54}{2}=27$\\
{\color{blue} --------------------}\\
{\bf (4a)}\\
$3x-10=2-x$\\
$3x+x=2+10$\\
$4x=12$\\
$x=\frac{12}{4}=3$\\
{\color{blue} --------------------}\\
{\bf (4b)}\\
$2(3p-6)=3p-9$\\
$6p-12=3p-9$\\
$6p-3p=-9+12$\\
$3p=3$\\
$p=\frac{3}{3}=1$\\
{\color{blue} --------------------}\\
{\bf (4c)}\\
$3(2y+1)-4(2y-3)=5$\\
$6y+3-8y+12=5$\\
$-2y+15=5$\\
$-2y=5-15$\\
$y=\frac{-10}{-2}=5$\\
\\
\\
\end{minipage}
\begin{minipage}{6cm}
{\color{blue} --------------------}\\
{\bf (4d)}\\
$\frac{1}{2}(x+3)=\frac{1}{4}(3x-1)$\\
$\left( \frac{1}{2}(x+3)=\frac{1}{4}(3x-1)\right)\times 8$\\
$4(x+3)=2(3x-1)$\\
$4x+12=6x-2$\\
$4x-6x=-2-12$\\
$-2x=-14$\\
$x=\frac{-14}{-2}=7$\\
{\color{blue} --------------------}\\
{\bf (5a)}\\
$\frac{2}{3}(5x+2)=2(x+2)$\\
$\left(\frac{2}{3}(5x+2)=2(x+2)\right)\times 3$\\
$2(5x+2)=3\cdot 2(x+2)$\\
$10x+4=6x+12$\\
$10x-6x=12-4$\\
$4x=8$\\
$x=\frac{8}{4}=2$\\
{\color{blue} --------------------}\\
{\bf (5b)}\\
$\frac{2y+4}{4}-\frac{4y-1}{3}=1$\\
$\left(\frac{2y+4}{4}-\frac{4y-1}{3}=1\right)\times 12$\\
$3(2y+4)-4(4y-1)=12$\\
$6y+12-16y+4=12$\\
$-10y=12-16$\\
$-10y=-4$\\
$y=\frac{-4}{-10}=\frac{2}{5}$\\
{\color{blue} --------------------}\\
{\bf (5c)}\\
$\frac{2x-1}{4}+\frac{6x+4}{2}=4x+3$\\
$\left(\frac{2x-1}{4}+\frac{6x+4}{2}=4x+3\right)\times 8$\\
$2(2x-1)+4(6x+4)=32x+24$\\
$4x-2+24x+16=32x+24$\\
$28x+14=32x+24$\\
$28x-32x=24-14$\\
$-4x=10$\\
$x=\frac{10}{-4}=-\frac{5}{2}$\\
\end{minipage}
\end{flushleft}
\end{document}
